\chapter{System calls and user pointers}
\label{ch:syscalls}

User programs cannot touch devices, other processes or the kernel's
memory; the CPU refuses. Everything they need from outside their own
address space, they must \emph{ask} for --- through a system call, a
controlled door into the kernel. Your \code{pthread\_create} is such a
door, and every argument that comes through it is chosen by the user
program: possibly wrong, possibly malicious.

\section{User mode and kernel mode}

x86 CPUs run in one of four privilege levels (rings); operating systems
use two. In \emph{kernel mode} (ring~0) every instruction is allowed. In
\emph{user mode} (ring~3) privileged instructions --- \code{cli},
loading \code{cr3}, I/O ports --- raise an exception, and pages whose
page-table entries lack the \emph{user} bit are inaccessible
(chapter~8). The only ways from user to kernel mode are:

\begin{description}
  \item[Interrupts] from devices (timer, keyboard) --- asynchronous, not
    caused by the running code;
  \item[Exceptions] caused by the current instruction: page fault,
    division by zero, general protection fault;
  \item[Traps] deliberate: the program executes \code{int~n} or
    \code{syscall} to request a service.
\end{description}

All three enter the kernel at an address the kernel registered in advance
(the \emph{interrupt descriptor table}, or the \code{syscall} MSRs). The
program chooses \emph{which} service --- by a number --- never \emph{where}
the kernel starts executing. An interrupt gate also has a privilege level
(DPL): user code may only trigger gates with DPL~3 via \code{int~n}.

\section{Calling conventions}

\begin{center}
\begin{tabular}{L{2.4cm}L{3.2cm}L{4.4cm}L{3.2cm}}
\toprule
 & \textbf{Instruction} & \textbf{Number / arguments} & \textbf{Result} \\
\midrule
Linux x86-64 & \code{syscall} & \code{rax} / \code{rdi, rsi, rdx, r10, r8, r9}
& \code{rax}, errors as $-$errno \\
SWEB x86-64 & \code{int \$0x80} & \code{rax} / \code{rbx, rcx, rdx, rsi, rdi}
& \code{rax} \\
\bottomrule
\end{tabular}
\end{center}

Linux uses \code{r10} for the fourth argument because the \code{syscall}
instruction itself overwrites \code{rcx} (return address) and \code{r11}
(flags). SWEB passes the number plus five arguments --- exactly enough for
\code{pthread\_create}'s four arguments plus the address of a start
wrapper.

\paragraph{Errors.} The Linux kernel returns $-$errno (a value between
$-4095$ and $-1$). The libc wrapper turns it into $-1$ and stores the
positive number in \code{errno}, a \emph{thread-local} variable
(chapter~2). Any other value, even a ``negative'' address, is a valid
result --- hence the test \code{(unsigned long)r > -4096UL}.

\section{One system call, end to end}

\begin{figure}[htbp]
\centering
\begin{tikzpicture}[font=\small,node distance=0.32cm,
    s/.style={draw,rounded corners=2pt,align=left,inner sep=4pt,text width=10.4cm},
    >=Latex]
  \node[s,fill=keybg] (a) {\code{write(1, "hi", 2)} \hfill \file{userspace/libc/src/write.c}};
  \node[s,fill=keybg,below=of a] (b) {\code{\_\_syscall(sc\_write, 1, buf, 2, 0, 0)}: \code{rax=4, rbx=1, rcx=buf, rdx=2};
    \code{int \$0x80} \hfill \file{arch/x86/64/userspace/syscalls.c}};
  \node[s,fill=black!8,below=of b] (c) {CPU: gate 0x80 (DPL 3), ring 3 $\to$ 0, switch to the thread's kernel stack,
    push user \code{rip, cs, rflags, rsp, ss}};
  \node[s,fill=newbg,below=of c] (d) {\code{arch\_syscallHandler}: save all user registers into
    \code{currentThread->user\_registers\_} \hfill \file{arch\_interrupts.S}};
  \node[s,fill=newbg,below=of d] (e) {\code{syscallHandler()}: interrupts on,
    \code{Syscall::syscallException(rax, rbx, rcx, rdx, rsi, rdi)} \hfill \file{InterruptUtils.cpp}};
  \node[s,fill=newbg,below=of e] (f) {\code{Syscall::write}: check the buffer, \code{kprintf} \hfill \file{Syscall.cpp}};
  \node[s,fill=newbg,below=of f] (g) {\code{user\_registers\_->rax = ret}; interrupts off; \code{arch\_contextSwitch()}
    restores the user registers, \code{iretq} to ring 3};
  \foreach \x/\y in {a/b,b/c,c/d,d/e,e/f,f/g} \draw[->] (\x) -- (\y);
  \draw[decorate,decoration={brace,amplitude=5pt}] ([xshift=4pt]a.north east) -- ([xshift=4pt]b.south east)
    node[midway,right=6pt] {user};
  \draw[decorate,decoration={brace,amplitude=5pt}] ([xshift=4pt]d.north east) -- ([xshift=4pt]g.south east)
    node[midway,right=6pt] {kernel};
\end{tikzpicture}
\caption{The path of \code{write} through SWEB (x86-64).}
\label{fig:syscallpath}
\end{figure}

Note two details in figure~\ref{fig:syscallpath}. The kernel runs the
system call on the \emph{thread's own kernel stack}, so every thread needs
one (in SWEB it is part of the \code{Thread} object). And interrupts are
\emph{on} during the call: a system call can be preempted, and several
threads can be inside the kernel at once.

\section{Never trust a user pointer}

For \code{write(fd, buf, len)} the kernel reads \code{len} bytes at
\code{buf}. If it did so blindly, a program could pass a kernel address
and have the kernel copy kernel memory into a file --- or an unmapped
address and crash the kernel. Every pointer from user space must be
checked:

\begin{enumerate}
  \item \textbf{The whole range is user memory.} Not just the start:
    \code{buf + len} must also stay below the user/kernel boundary.
  \item \textbf{The check cannot overflow.} \code{buf + len > USER\_BREAK}
    fails when \code{buf + len} wraps around: \code{buf = 0x1000},
    \code{len = 0xfffffffffffff800} gives \code{0x800}, ``below the
    break''. Never add two user-controlled numbers; subtract from a
    constant:
\begin{lstlisting}
if (len > USER_BREAK || buf > USER_BREAK - len) return -EFAULT;
\end{lstlisting}
  \item \textbf{The pages exist} (or the kernel handles the fault on a bad
    user page without dying). Linux copies with \code{copy\_from\_user},
    which turns a fault into $-$\code{EFAULT}; SWEB's page-fault handler
    loads valid user pages on demand and ends the process otherwise.
  \item \textbf{Output buffers are writable.}
  \item \textbf{Strings} have no known length: copy byte by byte (page by
    page) into a bounded kernel buffer, checking each page as the copy
    reaches it --- \code{strlen} on a user pointer \emph{is} the unchecked
    read.
\end{enumerate}

Not everything that looks like a pointer is one. \code{pthread\_create}'s
\code{arg} is never dereferenced by the kernel --- it is passed to the new
thread --- so it must \emph{not} be checked; NULL and 42 are legal.
\code{start\_routine} is not dereferenced either; it becomes an
instruction pointer and at most needs to lie below \code{USER\_BREAK}.

\section{TOCTOU and the double fetch}

Once a process has several threads, user memory can change \emph{while}
the kernel works with it. If the kernel reads a value from user memory to
check it and reads it again to use it, another thread can swap it in
between:

\begin{lstlisting}
if (!access_ok(iov[i].base, iov[i].len)) return -EFAULT;   /* fetch 1 */
copy_to_device(iov[i].base, iov[i].len);                   /* fetch 2: may differ! */
\end{lstlisting}

The fix is simple and absolute: copy the user data into kernel memory
\emph{once}, then check and use only the copy. (Modern CPUs add SMAP, which
makes the kernel fault on any user-memory access outside explicit
copy routines.)

\begin{swebbox}
Syscall numbers live in \file{common/include/kernel/syscall-definitions.h},
shared by the kernel and the libc. Adding a system call means: a number,
a \code{case} in \code{Syscall::syscallException}, a static
\code{Syscall} method, and a libc wrapper calling \code{\_\_syscall}
(\module{11}, part 5 does it for \code{gettid}). SWEB's own checks are
the hand-written \code{buffer + size > USER\_BREAK} kind --- including the
overflow hole above. Your \code{pthread\_create} and \code{pthread\_join}
receive user pointers (\code{pthread\_t *}, \code{void **}); check them
properly, and check them \emph{before} you create or free anything.
\end{swebbox}

\begin{keybox}
\begin{itemize}
  \item The user picks the service number, never the kernel entry point;
    gates and privilege levels enforce it.
  \item Linux: \code{syscall}, args in \code{rdi rsi rdx r10 r8 r9},
    $-$errno; SWEB: \code{int \$0x80}, number in \code{rax}, args in
    \code{rbx rcx rdx rsi rdi}.
  \item Every system call runs on the calling thread's kernel stack, with
    interrupts on.
  \item Validate the whole user range, overflow-free; copy strings with
    per-page checks; do not check opaque arguments.
  \item Fetch user data once; check and use the kernel copy.
\end{itemize}
\end{keybox}
